

\medskip
\noindent\SGARefTarget{sga3:iv:proposition:4:4:11}\textbf{Proposition 4.4.11.}
\textit{Let \(R\) be a \(\widetilde{\mathscr C}\)-equivalence relation on the
sheaf \(X\). For every subsheaf \(Y\) of \(X\) stable under \(R\), let
\(Y'=Y/R_Y\), regarded as a subsheaf of \(X'=X/R\). Then}
\[
  Y=Y'\times_{X'}X,
\]
\textit{and the maps}
\[
  Y\longmapsto Y/R_Y,
  \qquad
  Y'\longmapsto Y'\times_{X'}X
\]
\textit{define a bijective correspondence between the subsheaves \(Y\) of
\(X\) stable under \(R\) and the subsheaves \(Y'\) of \(X'\).}

If \(Y'\) is a subsheaf of \(X'\), then \(Y'\times_{X'}X\) is a subsheaf of
\(X\) stable under \(R\).\footnote{Translator's note: the Gille--Polo re-edition (2011), Expos\'e~IV, adds
that \((Y'\times_{X'}X)/R=Y'\), which gives the surjectivity of
\(Y\mapsto Y/R_Y\). Indeed \(Y'\times_{X'}X\to Y'\) is covering, as a base
change of the covering morphism \(X\to X'\) (4.4.9), so its sheaf image is
\(Y'\); by 4.4.10 this image is the quotient.}
If \(Y'\) is obtained by passage to the quotient from a subsheaf \(Y\) of
\(X\), then \(Y'\times_{X'}X\) contains \(Y\). It is therefore
enough to prove the following: if \(Y\) and \(Y_1\) are two subsheaves of
\(X\), both stable under \(R\), with \(Y_1\) containing \(Y\), and if their
quotients \(Y/R_Y\) and \(Y_1/R_{Y_1}\) are identical, then \(Y=Y_1\).

We may plainly reduce to the case \(Y_1=X\). Write
\[
  P=i(X)/i(R),\qquad Q=i(Y)/i(R_Y)
\]
for the corresponding presheaves. The diagram
\[
\begin{array}{ccc}
  X&\longrightarrow&P\\
  \big\uparrow&&\big\uparrow\\
  Y&\longrightarrow&Q
\end{array}
\]
is Cartesian. Since there is a commutative diagram
\[
\begin{array}{ccc}
  Q&\longrightarrow&a(Q)\\
  \big\downarrow&&\big\downarrow{\scriptstyle\sim}\\
  P&\longrightarrow&a(P).
\end{array}
\]
the monomorphism \(Q\hookrightarrow P\) is covering,\footnote{Translator's note: here \(Q\to a(Q)\) is covering (4.3.11),
and \(a(Q)\to a(P)\) is an isomorphism, the two quotients being identical; as
\(P\) is separated (4.4.5 (ii)), \(P\to a(P)\) is a monomorphism, so that
\(Q\times_PS=Q\times_{a(P)}S\) for every \(S\to P\). The Gille--Polo re-edition
(2011), Expos\'e~IV, cites 4.3.11 here.} so \(Q\) is a refinement of \(P\).
After base change, \(Y\) is a refinement of \(X\). Because \(X\) and \(Y\)
are sheaves, \SGARefLink{sga3:iv:corollary:4:3:12}{4.3.12} then gives \(Y=X\).

\medskip
\noindent\SGARefTarget{sga3:iv:para:4:4:12}\textbf{4.4.12.}
In particular, if \(Y\) is a subsheaf of \(X\) and \(Y'=Y/R_Y\), the
preceding correspondence defines a subsheaf \(\overline Y\) of \(X\), stable
under \(R\), containing \(Y\), and minimal with these properties.
It is called the \emph{saturation} of \(Y\) for the equivalence relation
\(R\).
